\begin{tikzpicture}[font=\small,>=stealth]
\draw[gray!45,rounded corners=3pt] (0,0) rectangle (14.2,9.5);
\node at (7.1,8.85) {方法二：两腰按份数画出，差定就能求和};

\node at (7.1,7.9) {先取 $\lambda=2$：长腰两份、短腰一份。};
\draw[orange!85!black,line width=4pt] (1.1,6.7)--(9.1,6.7);
\foreach \x in {1.1,5.1,9.1}\draw (\x,6.5)--(\x,6.9);
\node[anchor=east] at (.9,6.7) {$r_1$};
\draw[purple!80!black,line width=4pt] (1.1,4.6)--(5.1,4.6);
\foreach \x in {1.1,5.1}\draw (\x,4.4)--(\x,4.8);
\node[anchor=east] at (.9,4.6) {$r_2$};
\draw[<->] (5.1,5.9)--(9.1,5.9) node[midway,below=4pt] {差：$1$ 份 $=2a$};
\node[anchor=west] at (6.1,4.6) {所以一份长 $2a$};
\draw[orange!85!black,line width=4pt] (1.1,2.5)--(9.1,2.5);
\draw[purple!80!black,line width=4pt] (9.1,2.5)--(13.1,2.5);
\node at (7.1,1.7) {和：$3$ 份 $=6a$，必须不小于 $2c$};
\node at (7.1,.65) {一般情形：差 $|\lambda-1|$ 份、和 $\lambda+1$ 份。};

\end{tikzpicture}
